\section{Introduction}\label{sec:intro}

Multi-Component Multi-Layer Neural Networks (MMNNs) or RF-LR combine frozen random features with low-rank bottlenecks, reducing the parameter budget from \( O(N^2) \) to \( O(rN) \) while retaining NTK-regime convergence~\cite{zhang2025structuredbalancedmulticomponentmultilayer}. We develop an NTK theory for RF-LR that proves: (i)~the same RKHS as the shallow ReLU kernel (no expressivity loss); (ii)~rank-driven concentration with sub-Gaussian tails in \( r \); (iii)~explicit Puiseux expansions via Fisher--Kibble decoupling and hypergeometric analysis.

\paragraph{Contributions.} We derive the NTK recursion with \( 1/r \) scaling at bottlenecks (Section~\ref{sec:network-definition}), prove the mean NTK induces the same RKHS as the shallow kernel (Section~\ref{sec:no-rkhs-advantage}) via a \textbf{hypergeometric proof in the main text} that computes \( \mathbb{E}[\arccos(\hat{\rho}_r)] \) over Fisher's distribution, and link the \( I(r) \sim 1/\sqrt{r} \) decay to tighter concentration (Proposition~\ref{prop:I-r-scaling}). We state the spike--Marchenko--Pastur spectral structure (Section~\ref{sec:spectra}).

\paragraph{Related work.} NTK theory~\cite{jacot2018neural} has been extended to finite width~\cite{arora2019exact} and depth~\cite{terjek2022depth}. The ``deep equals shallow'' RKHS result~\cite{bietti2021deepequalsshallowrelu} characterizes spectral decay via endpoint expansions. We extend this to RF-LR using Fisher--Kibble decoupling~\cite{fisher1915probable,hotelling1953new}, avoiding correlation-diagram computations.
